\documentclass[11pt,a4paper]{article}
\usepackage[T1]{fontenc}
\usepackage{lmodern}
\usepackage[margin=1in]{geometry}
\usepackage{amsmath,amssymb,amsthm,mathtools}
\usepackage{microtype}
\usepackage[colorlinks=true,linkcolor=blue,citecolor=blue,urlcolor=blue]{hyperref}
\newtheorem{theorem}{Theorem}[section]
\newtheorem{lemma}[theorem]{Lemma}
\newcommand{\R}{\mathbb R}
\newcommand{\N}{\mathbb N}
\newcommand{\dd}{\,\mathrm d}
\newcommand{\supp}{\operatorname{supp}}
\newcommand{\spec}{\sigma}
\title{A smooth counterexample to critical--subcritical heat-kernel comparison}
\author{Kaiwen Jin\\\small\href{mailto:yui.ui.siki@gmail.com}{yui.ui.siki@gmail.com}}
\date{October 15, 2026}
\begin{document}
\maketitle
\begin{abstract}
We construct two nonnegative Schr\"odinger operators on the line with smooth
bounded potentials and spectral bottom zero. One is subcritical and the other
null-critical, while the quotient of their diagonal heat kernels has infinite
limsup. This gives a negative answer to the heat-kernel comparison question
without an ordering assumption on the potentials. The construction preserves
small heat-kernel values of successive operators with positive spectral gaps,
then restores criticality by adding distant intervals on which a positive
zero-energy solution equals one.
\end{abstract}

\section{The comparison question}

The distinction between critical and subcritical Schr\"odinger operators can be
read from the time integral of the positive heat kernel. It is natural to ask
whether this distinction also orders their decay at large times. Fraas,
Krej\v{c}i\v{r}\'ik and Pinchover formulated such a comparison in
\cite[Conjecture 1]{FKP}; its self-adjoint heat-kernel version appears again in
\cite[Conjecture 2]{Krejcirik}. Their conditional results include an additional
global heat-kernel comparison hypothesis \cite[Theorem 3.1]{FKP} and
nonnegative potential perturbations \cite[Corollary 1]{FKP}.

We consider the formulation in AIM Problem 217 \cite{AIM217}. On a connected,
complete, noncompact smooth Riemannian manifold without boundary, let
$H_j=-\Delta_g+V_j$, $j\in\{0,+\}$, be the Friedrichs operators for smooth
real potentials bounded below. Assume $H_j\geq0$ and $\inf\spec(H_j)=0$.
Write $k_j(t,x,y)>0$ for the heat kernel of $e^{-tH_j}$. At some fixed
distinct points $x_*,y_*$, suppose that
\begin{equation}\label{eq:original-integrals}
 \int_0^\infty k_+(t,x_*,y_*)\dd t<\infty,
 \qquad
 \int_0^\infty k_0(t,x_*,y_*)\dd t=\infty,
 \qquad \ker_{L^2}H_0=\{0\}.
\end{equation}
These are the subcritical and null-critical cases, respectively. With no
pointwise ordering of $V_+$ and $V_0$, the question is whether
\begin{equation}\label{eq:question}
 \lim_{t\to\infty}\sup_{(x,y)\in K}
       \frac{k_+(t,x,y)}{k_0(t,x,y)}=0
 \quad\text{for every compact }K\subset M\times M.
\end{equation}

\begin{theorem}\label{thm:main}
There are real potentials $V_+,V_0\in C^\infty(\R)\cap L^\infty(\R)$ such
that their Friedrichs operators $H_j=-\mathrm d^2/\mathrm dx^2+V_j$ on
$L^2(\R,\mathrm dx)$ satisfy
\[
 H_j\geq0,\qquad \inf\spec(H_j)=0,\qquad \ker_{L^2}H_0=\{0\},
\]
and their positive heat kernels satisfy
\begin{equation}\label{eq:main-integrals}
 \int_0^\infty k_+(t,0,1)\dd t<\infty,
 \qquad \int_0^\infty k_0(t,0,1)\dd t=\infty,
\end{equation}
but
\begin{equation}\label{eq:main-ratio}
 \limsup_{t\to\infty}\frac{k_+(t,0,0)}{k_0(t,0,0)}=\infty.
\end{equation}
\end{theorem}

The Euclidean line is an admissible manifold. The compact singleton
$K=\{(0,0)\}$ therefore disproves \eqref{eq:question}; the distinctness of
the points in \eqref{eq:original-integrals} imposes no restriction on this $K$.

The numerator will be a fixed nonzero nonnegative compactly supported bump.
Its diagonal heat kernel is bounded below by $ct^{-3/2}$. For the reference
operator we use a positive zero-energy solution $h=e^w$ and set
$V_0=w''+(w')^2$. A finite construction with $w'=1$ on its positive tail has
a strictly positive spectral gap, so its heat kernel is eventually much
smaller than $ct^{-3/2}$. We choose one such time, then a radius so large
that later changes outside that radius preserve the estimate. Beyond the
radius we let $w$ return to zero on a unit interval. Repeating this produces
one fixed potential with infinitely many such intervals. They make the
reference Green integral infinite without losing the selected small kernel
values. The analytic estimates are given in Section~\ref{sec:estimates};
Section~\ref{sec:construction} carries out the construction. The
one-dimensional Dirichlet and Green-kernel justification is proved in
Appendix~\ref{sec:dirichlet}.

\section{Heat-kernel estimates}\label{sec:estimates}

For a bounded real continuous potential $U$, write
$H_U=-\mathrm d^2/\mathrm dx^2+U$. Its self-adjoint domain is $H^2(\R)$,
its form domain is $H^1(\R)$, and its Friedrichs form is
\begin{equation}\label{eq:form}
 Q_U[u]=\int_\R \bigl(|u'|^2+U|u|^2\bigr)\dd x.
\end{equation}
These identifications follow from the free operator and bounded perturbation
\cite[Theorems 6.4 and 7.8, Lemma 7.9]{Teschl}. One can check the form core
directly: first multiply an $H^1$ function by smooth cutoffs whose derivatives
are $O(R^{-1})$, then mollify. Both the function and its weak derivative
converge in $L^2$, so $C_c^\infty(\R)$ is a form core.

Our free kernel and Brownian-bridge convention are
\begin{equation}\label{eq:fk}
 \begin{split}
 p(t,x,y)&=(4\pi t)^{-1/2}e^{-(x-y)^2/(4t)},\\
 k_U(t,x,y)&=p(t,x,y)\,
  \mathbb E\exp\left(-t\int_0^1 U(b_r)\dd r\right),\\
 b_r&=(1-r)x+ry+\sqrt{2t}\,\beta_r,
 \end{split}
\end{equation}
where $\beta$ is the standard unit-time Brownian bridge. This is the
Feynman--Kac formula \cite[Theorem 1.1]{Simon}, applied to $U/2$ at time
$2t$. In particular, $k_U$ is strictly positive and continuous for $t>0$.
If $U\geq-B$, then $k_U\leq e^{Bt}p$. The symmetry and semigroup identities
hold for these continuous representatives, and their convolution integrals
are controlled by the Gaussian bound.

\begin{lemma}\label{lem:locality}
If bounded continuous real potentials $U,W\geq-B$ agree on $[-R,R]$, then
\begin{equation}\label{eq:locality}
 |k_U(t,0,0)-k_W(t,0,0)|
 \leq 4e^{Bt}(4\pi t)^{-1/2}e^{-R^2/t},\qquad t>0.
\end{equation}
\end{lemma}
\begin{proof}
Reflection of Brownian paths at their first hit of a level preserves Wiener
law \cite[Proposition 4]{Lalley}. For the variance-two bridge from $x$ to
$y$, with $x,y<a$, division of the reflected endpoint density by the free
endpoint density gives
\begin{equation}\label{eq:reflection}
 \mathbb P\{\max b_r\geq a\}
   =\frac{p(t,x,2a-y)}{p(t,x,y)}
   =e^{-(a-x)(a-y)/t}.
\end{equation}
The lower-level formula follows by negating the path. A bridge from zero to
zero thus exits $[-R,R]$ with probability at most $2e^{-R^2/t}$. On paths
that do not exit, the two weights in \eqref{eq:fk} coincide; on the remaining
paths their difference has absolute value at most $2e^{Bt}$. Multiplying
these bounds by $p(t,0,0)$ proves \eqref{eq:locality}.
\end{proof}

We also need an exact formula for the time-integrated kernel. We state it
here and prove the Dirichlet identification and exhaustion in
Appendix~\ref{sec:dirichlet}, so that the construction can be read separately
from that operator argument.

\begin{lemma}\label{lem:green}
Let $U$ be smooth and bounded on $\R$, and let $h>0$ solve $-h''+Uh=0$.
Put $s(x)=\int_0^x h(r)^{-2}\dd r$. For $a>\max(|x|,|y|)$, define
\begin{equation}\label{eq:green}
 G_a(x,y)=h(x)h(y)
 \frac{[s(\min(x,y))-s(-a)]
       [s(a)-s(\max(x,y))]}{s(a)-s(-a)}.
\end{equation}
Then
\begin{equation}\label{eq:green-exhaustion}
 \int_0^\infty k_U(t,x,y)\dd t
    =\lim_{a\to\infty}G_a(x,y),
\end{equation}
where the limit may be infinite and may be taken along increasing integers.
If $D=(a,\infty)$ and $U=0$ on $D$, its Dirichlet kernel is
\begin{equation}\label{eq:image}
 p_D(t,u,v)=p(t,u,v)-p(t,u,2a-v),\qquad u,v>a,
\end{equation}
and $0\leq p_D(t,u,v)\leq k_U(t,u,v)$.
\end{lemma}

\begin{lemma}\label{lem:numerator}
Set
\[
 q(x)=\begin{cases}e^{-1/(1-x^2)},&|x|<1,\\0,&|x|\geq1.\end{cases}
\]
For $H_+=H_q$, one has $H_+\geq0$, $\inf\spec(H_+)=0$ and
$\int_0^\infty k_+(t,0,1)\dd t<\infty$. There is a fixed $c>0$ such that
\begin{equation}\label{eq:numerator-lower}
 k_+(t,0,0)\geq ct^{-3/2}\qquad(t\geq3).
\end{equation}
\end{lemma}
\begin{proof}
The bump is smooth, even and nonnegative, so $H_+\geq0$. Take
$\psi\in C_c^\infty((1,2))$ with $\|\psi\|_2=1$ and set
$\psi_L(x)=L^{-1/2}\psi((x-2)/L)$ for $L\geq1$. These functions are
supported where $q=0$ and have Rayleigh quotient
$L^{-2}\|\psi'\|_2^2$. Hence $\inf\spec(H_+)=0$.

Let $h_+''=qh_+$, with $h_+(0)=1$, $h_+'(0)=0$. For completeness, the
Volterra equation
\[
 h_+(x)=1+\int_0^x(x-r)q(r)h_+(r)\dd r\quad(x\geq0)
\]
constructs this solution: its $m$th iterated term is bounded by
$\|q\|_\infty^m x^{2m}/(2m)!$ on compact intervals. The series and its
differentiated series converge there, and the same bound for a difference
gives uniqueness. All terms are nonnegative. Thus $h_+\geq1$,
$h_+'\geq0$, and $b=h_+'(1)>0$. Even reflection gives the global smooth
positive solution. If $A=h_+(1)$, then for $x\geq1$,
\[
 h_+(x)=A+b(x-1),\qquad
 S_+:=\int_0^\infty h_+^{-2}
     =\int_0^1 h_+^{-2}+\frac1{bA}<\infty.
\]
In \eqref{eq:green}, $s$ is odd and tends to $\pm S_+$ at the two ends.
Every factor at $(x,y)=(0,1)$ has a finite limit, and the denominator tends
to $2S_+>0$. Equation~\eqref{eq:green-exhaustion} proves the finite integral.

To obtain a lower bound, let $I=[2,3]$, $D=(1,\infty)$ and
$m=\min_{u\in I}k_+(1,0,u)>0$. The semigroup identity, positivity, and
Lemma~\ref{lem:green} give, for $s=t-2\geq1$,
\[
 k_+(t,0,0)\geq m^2\int_I\int_I p_D(s,u,v)\dd u\dd v.
\]
On $I^2$, the image formula becomes
\[
 p_D(s,u,v)=(4\pi s)^{-1/2}e^{-(u-v)^2/(4s)}
                 \bigl[1-e^{-(u-1)(v-1)/s}\bigr].
\]
Here $|u-v|\leq1$, $(u-1)(v-1)\geq1$, and concavity gives
$1-e^{-1/s}\geq(1-e^{-1})/s$. Since $s\leq t$ and $I^2$ has area one,
\eqref{eq:numerator-lower} follows with
\begin{equation}\label{eq:constant-c}
 c=\frac{m^2e^{-1/4}(1-e^{-1})}{\sqrt{4\pi}}>0.
\end{equation}
\end{proof}

\section{Construction of the reference potential}\label{sec:construction}

We use one fixed smooth step
\begin{equation}\label{eq:theta}
 \theta(r)=\frac{e^{-1/r}}{e^{-1/r}+e^{-1/(1-r)}}\quad(0<r<1),
 \qquad \theta=0\ (r\leq0),\quad \theta=1\ (r\geq1).
\end{equation}
It is flat at both endpoints, strictly increasing on $(0,1)$, and satisfies
$\theta(1-r)=1-\theta(r)$. In particular,
$\int_0^1\theta=1/2$ and $B:=\|\theta'\|_\infty<\infty$. Consider smooth,
even, nonnegative functions $w$ with
\begin{equation}\label{eq:w-bounds}
 |w'|\leq1,\qquad |w''|\leq B.
\end{equation}
For each such function set $U_w=w''+(w')^2$ and $h=e^w$. Direct integration
by parts on compactly supported smooth functions gives
\begin{equation}\label{eq:factorization}
 Q_{U_w}[u]=\|u'-w'u\|_2^2\geq0.
\end{equation}
Indeed, the cross term is $-\int w'(|u|^2)'=\int w''|u|^2$, also for
complex $u$. The bounds \eqref{eq:w-bounds} and the form-core approximation
extend \eqref{eq:factorization} to $H^1(\R)$. Thus $H_{U_w}\geq0$ and
$-h''+U_wh=0$. Any $L^2$ zero mode belongs to $H^2$ and satisfies
$u'-w'u=0$, whence $(u/h)'=0$ on each compact interval and $u=Ch$.
In particular, $h\notin L^2$ excludes every nonzero $L^2$ zero mode.

\begin{lemma}\label{lem:gap}
If, in addition, $w'=1$ for all sufficiently large positive $x$, then
$\lambda(w):=\inf\spec(H_{U_w})>0$. For $t\geq1$,
\begin{equation}\label{eq:gap-kernel}
 k_{U_w}(t,0,0)\leq M e^{-\lambda(w)(t-1)},
 \qquad M=\frac{e^B}{\sqrt{4\pi}}.
\end{equation}
\end{lemma}
\begin{proof}
Evenness implies $U_w=1$ outside a compact interval, and $e^w$ grows
exponentially at both ends. If the spectral bottom were zero, the spectral
form characterization \cite[Theorem 3.7 and equations (3.50)--(3.51)]{Teschl}
would give $u_\ell\in H^1$ with $\|u_\ell\|_2=1$ and
$Q_{U_w}[u_\ell]\to0$. Since $U_w\geq-B$, this sequence is bounded in
$H^1$. Pass to a weakly convergent subsequence there
\cite[Lemma 1.13]{Teschl}, with limit $u$, and to a further subsequence
converging strongly in $L^2$ on every compact interval. The last assertion
can be checked in one dimension as follows. On any fixed interval,
\[
 |v(x)-v(y)|\leq\|v'\|_2|x-y|^{1/2};
\]
together with a point where $|v|$ is at most its $L^2$ average, this gives
uniform bounds and a common modulus of continuity. Successively finer finite
grids give a subsequence converging uniformly on that interval. A diagonal
choice over $[-j,j]$, $j\in\N$, proves the stated local convergence.

If $u=0$, compact support and boundedness of $U_w-1$ would give
\[
 Q_{U_w}[u_\ell]=\|u_\ell'\|_2^2+1+
                 \int(U_w-1)|u_\ell|^2\geq1+o(1),
\]
a contradiction. Thus $u\ne0$. But \eqref{eq:factorization} implies
$u_\ell'-w'u_\ell\to0$ strongly in $L^2$, so its weak limit is
$u'-w'u=0$. This forces $u=Ce^w$, impossible for a nonzero $L^2$ function.
The gap is therefore strictly positive.

For $f(z)=k_{U_w}(1/2,z,0)$, the Gaussian bound gives $f\in L^2$ and
$\|f\|_2^2=k_{U_w}(1,0,0)$. The semigroup identity and spectral bound give
\[
 k_{U_w}(t,0,0)\leq e^{-\lambda(w)(t-1)}\|f\|_2^2
       \leq e^{-\lambda(w)(t-1)}e^B(4\pi)^{-1/2}.
\]
The constant $M$ is common to these functions; their positive gaps need not
have a common lower bound.
\end{proof}

The next extension lets a function with an outward slope return to height
zero. Suppose $R$ lies on the slope-one tail and $a=w(R)\geq1/2$. Keep $w$
on $[-R,R]$ and prescribe its derivative on the positive side as follows:
\begin{enumerate}
\item On $[R,R+2]$, use $1-2\theta((x-R)/2)$. The slope turns from one
to minus one, with total height change zero; throughout this turn the height
is at least $a$.
\item Continue with slope $-1$ for length $a-1/2$, reaching height $1/2$.
\item On the next unit interval, with local coordinate $r\in[0,1]$, use
derivative $-1+\theta(r)$. The height decreases by $1/2$ and reaches zero.
\item Keep height zero on a unit interval $J$.
\item On the next unit interval use derivative $\theta(r)$, reaching
height $1/2$ with slope one, and continue with slope one thereafter.
\end{enumerate}
Every piece stays nonnegative. Flatness of $\theta$ matches all jets, even
when the straight descent has zero length. The second derivative in the
first turn is $-\theta'((x-R)/2)$ and in each unit turn has absolute value
at most $B$. Thus \eqref{eq:w-bounds} is preserved. Reflect on the negative
side; the central part is unchanged. Denote this extension by $w_{\rm new}$.

\begin{proof}[Proof of Theorem~\ref{thm:main}]
Fix $V_+=q$ and the constant $c$ of Lemma~\ref{lem:numerator}. Start from
\[
 w_1(x)=\int_0^{|x|}\theta(r-1)\dd r.
\]
This function is smooth and zero on $[-1,1]$, obeys
\eqref{eq:w-bounds}, and has positive-side slope one for $x\geq2$.
Set $t_0=R_0=0$.

Given $w_n$, let $U_n=w_n''+(w_n')^2$ and
$\lambda_n=\inf\spec(H_{U_n})>0$, using Lemma~\ref{lem:gap}. First choose
an integer $t_n$ such that
\begin{equation}\label{eq:time-choice}
 t_n\geq\max(n+3,t_{n-1}+1),\qquad
 M e^{-\lambda_n(t_n-1)}\leq\frac{c}{2n}t_n^{-3/2}.
\end{equation}
This is possible because $\lambda_n$ is positive for this fixed $w_n$.
Next choose an integer $R_n$ on its slope-one tail, beyond every earlier
valley, with $R_n>\max(n,R_{n-1}+1)$ and $w_n(R_n)\geq1$, such that
\begin{equation}\label{eq:radius-choice}
 4e^{Bt_n}(4\pi t_n)^{-1/2}e^{-R_n^2/t_n}
       \leq\frac{c}{2n}t_n^{-3/2}.
\end{equation}
Take $R_n$ at least one unit beyond the beginning of that tail as well.
All these conditions involve finitely many fixed quantities, and the left
side of \eqref{eq:radius-choice} tends to zero as $R_n$ increases.
Apply the extension just described at $R_n$ to obtain $w_{n+1}$ and its new
unit zero interval $J_n$. If desired, taking the first admissible integers
at both choices makes the construction fully specified without a numerical
evaluation of $\lambda_n$.

We have $w_{n+1}=w_n$ on $[-R_n,R_n]$, with $R_n\to\infty$. Every compact
neighborhood is therefore eventually unchanged, with all derivatives. Its
eventual value defines one smooth even nonnegative function $w$ obeying
\eqref{eq:w-bounds}. It is zero on $[-1,1]$ and on all the disjoint unit
intervals $J_n$ and $-J_n$. Define the single final potential and solution
\begin{equation}\label{eq:final-potential}
 V_0=w''+(w')^2,\qquad h_0=e^w.
\end{equation}
Then $-B\leq V_0\leq B+1$, $H_0\geq0$ by
\eqref{eq:factorization}, and $\ker_{L^2}H_0=\{0\}$ because $h_0\geq1$.

The preserved unit zero intervals give
\begin{equation}\label{eq:infinite-resistance}
 \int_0^\infty h_0^{-2}=\int_{-\infty}^0 h_0^{-2}=\infty.
\end{equation}
Put $s(x)=\int_0^x h_0^{-2}$ and $S=s(a)$. Evenness makes $s$ odd, and
$S\to\infty$ as $a\to\infty$. On the central interval $h_0=1$, so
$s(1)=1$. For $a>1$, \eqref{eq:green} at $(0,1)$ reads
\[
 G_a(0,1)=\frac{S-1}{2}.
\]
Lemma~\ref{lem:green} proves the infinite reference integral in
\eqref{eq:main-integrals}. The finite numerator integral has already been
proved at the same distinct pair.

For the spectral bottom, use the compactly supported $H^1$ function
\[
 u_a(x)=h_0(x)\max\left(0,1-\frac{|s(x)|}{S}\right).
\]
It vanishes outside $[-a,a]$, and its zero endpoint values ensure that its
zero extension is in $H^1$. Equation~\eqref{eq:factorization} yields
\[
 Q_0[u_a]=\frac1{S^2}\int_{-a}^a h_0^{-2}\dd x=\frac2S.
\]
If $S\geq2$, its square is at least $1/4$ on $[-1,1]$, whence
$\|u_a\|_2^2\geq1/2$. The Rayleigh quotients are at most $4/S\to0$;
together with nonnegativity, this gives $\inf\spec(H_0)=0$.

Finally, $V_0=U_n$ on $[-R_n,R_n]$, and both have the same lower bound
$-B$. Lemmas~\ref{lem:locality} and~\ref{lem:gap}, followed by
\eqref{eq:time-choice}--\eqref{eq:radius-choice}, give
\[
 0<k_0(t_n,0,0)\leq k_{U_n}(t_n,0,0)+\frac{c}{2n}t_n^{-3/2}
                  \leq\frac cn t_n^{-3/2}.
\]
The same fixed $c$ appears in \eqref{eq:numerator-lower}. Consequently
\[
 \frac{k_+(t_n,0,0)}{k_0(t_n,0,0)}\geq n,
\]
which proves \eqref{eq:main-ratio}. Both potentials are fixed; only the
comparison times tend to infinity.
\end{proof}

The constructed potentials are unordered. At the middle of each length-two
downward turn, $w'=0$ and $w''=-\theta'(1/2)<0$, so $V_0<0$ outside the
support of $q$. On a retained slope-one segment $V_0=1$ while $q=0$.
Thus $q-V_0$ takes both signs, as is consistent with the nonnegative
perturbation result in \cite[Corollary 1]{FKP}.

\appendix
\section{Dirichlet kernels and Green functions}\label{sec:dirichlet}

We prove Lemma~\ref{lem:green}, including the identification of killed paths
with the one-dimensional Friedrichs Dirichlet operator. Only the whole-line
formula \eqref{eq:fk} is imported from \cite{Simon}. The form, weak compactness,
spectral calculus and uniform polynomial approximation used below are given
in \cite[Theorems 2.12--2.13, Lemma 1.13, Theorems 3.1, 3.7 and 1.15]{Teschl}.

\subsection*{The killed semigroup}
Let $I$ be a finite open interval or a half-line, and let $U\geq-B$ be
bounded, real and continuous on $\R$. Define $H_U^I$ on $L^2(I)$ by the
closed form $Q_I[u]=\int_I(|u'|^2+U|u|^2)$ with domain $H^1_0(I)$.
Set
\[
 F(x)=\min\{1,\operatorname{dist}(x,\overline I)^2\},\quad
 H_m=H_{U+mF},\quad \alpha=B+1.
\]
The shifted forms
\[
 a_m[u]=\int_\R\bigl(|u'|^2+(U+\alpha)|u|^2+mF|u|^2\bigr)
          \geq\|u\|_{H^1}^2
\]
are coercive. For $f\in L^2(\R)$, put $u_m=(H_m+\alpha)^{-1}f$.
Testing its variational equation with $u_m$ gives
\[
 \|u_m\|_{H^1}\leq\|f\|_2,\qquad
 m\int F|u_m|^2\leq\|f\|_2^2.
\]
A weak $H^1$ subsequential limit $u$ vanishes outside $\overline I$, since
on every compact set where $F\geq\delta>0$ the local $L^2$ norm tends to
zero. An $H^1$ function in one dimension is locally absolutely continuous:
integrate its weak derivative and subtract, leaving a distribution with
derivative zero. Thus $u$ has zero trace at every finite endpoint of $I$.

These zero extensions are precisely $H^1_0(I)$. To check the needed density
at an endpoint $a$, use $u(a)=0$ and
\[
 |u(x)|^2\leq |x-a|\int_{\min(a,x)}^{\max(a,x)}|u'|^2.
\]
A cutoff vanishing within distance $\delta$ of $a$ and equal to one beyond
$2\delta$ has derivative-of-cutoff error bounded by a constant times
$\int_{|x-a|<2\delta}|u'|^2$, which tends to zero. Ordinary cutoffs at
infinity, followed by mollification away from the endpoints, complete the
$H^1$ approximation. Translation continuity in $L^2$ justifies convergence
for both $u$ and $u'$; equivalently this follows from the Fourier and weak
derivative description \cite[Section 7.1, (7.12)--(7.15)]{Teschl}.

Tests compactly supported in $I$ remove the term $mF$ from the variational
equation. Passing to the weak limit and using this density identifies
$u=(H_U^I+\alpha)^{-1}(f|_I)$, extended by zero. The solution is unique, so
every weak subsequential limit is the same. Since $Fu=0$, the variational
equations give the exact identity
\[
 a_m[u_m-u]=\operatorname{Re}\langle f,u_m-u\rangle\longrightarrow0.
\]
Hence $u_m\to u$ strongly in $H^1$ and $L^2$. In particular, the bounded
positive resolvents $R_m=(H_m+\alpha)^{-1}$ converge strongly to the
Dirichlet resolvent on $L^2(I)$, extended by zero on $L^2(I^c)$; call this
limit $R_\infty$.

The limiting resolvent vanishes on $L^2(I^c)$, so we pass to semigroups
through bounded continuous calculus. For fixed $t>0$, the
function
\[
 f_t(r)=e^{-t(1/r-\alpha)}\quad(r>0),\qquad f_t(0)=0
\]
is continuous on $[0,1]$. The spectra of $R_m$ and $R_\infty$ lie in that
interval. Strong convergence and uniform boundedness give convergence of
every finite power. Uniform polynomial approximation then gives
$f_t(R_m)\to f_t(R_\infty)$ strongly. The spectral calculus identifies this
as convergence of $e^{-tH_m}$ to $e^{-tH_U^I}$ extended by zero.

For the passage from the penalized kernels, recall the variance-two bridge
and its free prefactor:
\[
 b_r=(1-r)x+ry+\sqrt{2t}\,\beta_r,\qquad
 p(t,x,y)=(4\pi t)^{-1/2}e^{-(x-y)^2/(4t)},
\]
where $\beta$ is the standard unit-time bridge. In the full-line formula
\eqref{eq:fk} for $U+mF$, the weights decrease to
\[
 \exp\left(-t\int_0^1 U(b_r)\dd r\right)
       \mathbf1_{\{b_r\in\overline I\ \text{for all }r\}}.
\]
Continuity of each path identifies the zero integral of $F$ with this
event. For an upper level $z>\max(x,y)$, the reflection formula
\eqref{eq:reflection} reads
\[
 \mathbb P\{\max_{0\leq r\leq1}b_r\geq z\}
       =e^{-(z-x)(z-y)/t}.
\]
Its continuity in $z$ excludes atoms of the maximum above both endpoints;
negating the path gives the corresponding statement for the minimum.
Thus replacing $\overline I$ by $I$ changes a null
event for interior endpoints. The Gaussian bound allows passage to the
limit against compactly supported $L^2$ tests, identifying the strong
semigroup limit. We obtain the killed kernel
\begin{equation}\label{eq:killed}
 k_U^I(t,x,y)=p(t,x,y)\,
  \mathbb E\left[
   e^{-t\int_0^1U(b_r)\dd r}
   \mathbf1_{\{b_r\in I\ \text{for all }r\}}\right].
\end{equation}
This identity is pointwise: coupling bridges with the same $\beta$ proves
continuity at interior $(t,x,y)$ by dominated convergence, since failure of
convergence of the survival indicator requires a boundary maximum or
minimum. In particular, $0\leq k_U^I\leq k_U$ and $k_U^I\leq e^{Bt}p$.
For $I=(a,\infty)$ and $U=0$ on $I$, reflection gives the image formula
\eqref{eq:image}. For $I=(-a,a)$ increasing to $\R$, every bridge has
bounded range, and its survival indicators increase to one. Hence
\begin{equation}\label{eq:heat-exhaustion}
 k_U^{(-a,a)}(t,x,y)\uparrow k_U(t,x,y)
 \quad\text{at each fixed }t>0,x,y.
\end{equation}

\subsection*{The interval inverse and its time integral}
Now suppose $U$ is smooth and bounded, and $h>0$ solves $-h''+Uh=0$.
On $(-a,a)$ put $m_a=\min h>0$, $M_a=\max h$, and
$s(x)=\int_0^xh^{-2}$. Multiplication and division by $h$ preserve
$H^1_0(-a,a)$. For $u=hv$, integration by parts and density give
\[
 Q_I[u]=\int_Ih^2|v'|^2
 \geq\frac{m_a^2}{4a^2M_a^2}\|u\|_2^2.
\]
Here $\|v\|_2^2\leq4a^2\|v'\|_2^2$ follows by integrating
$v(x)=\int_{-a}^xv'$. Thus the Dirichlet operator has a positive gap
$\lambda_a=m_a^2/(4a^2M_a^2)$, allowed to depend on $a$.

The homogeneous solutions
$\ell(x)=h(x)[s(x)-s(-a)]$ and $r(x)=h(x)[s(a)-s(x)]$ satisfy
\[
 \ell'r-\ell r'=s(a)-s(-a)>0.
\]
Their Green function is
\[
 G_a(x,y)=\frac{\ell(\min(x,y))r(\max(x,y))}{s(a)-s(-a)},
\]
which is exactly \eqref{eq:green}. It is continuous, zero at the
endpoints, and its $x$ derivative has jump $-1$ at $x=y$. Thus
$H_U^I G_a(\cdot,y)=\delta_y$. Integration against a smooth compactly
supported function gives the Dirichlet inverse, whose uniqueness follows
from the positive gap.

Spectral integration on $[\lambda_a,\infty)$ gives
$\int_0^\infty e^{-tH_U^I}\dd t=(H_U^I)^{-1}$. The kernel equality also
holds at every point. For $0<t\leq1$, \eqref{eq:killed} is bounded by
$e^B(4\pi t)^{-1/2}$, integrable at zero in dimension one. For $t\geq1$,
the time-$1/2$ kernel vectors and the spectral bound give
\[
 |k_U^I(t,x,y)|
 \leq e^{-\lambda_a(t-1)}
       \sqrt{k_U^I(1,x,x)k_U^I(1,y,y)}
 \leq\frac{e^B}{\sqrt{4\pi}}e^{-\lambda_a(t-1)}.
\]
These bounds make the time-integrated kernel continuous, including on the
diagonal. Its equality with the inverse kernel, initially almost everywhere,
therefore holds everywhere. Finally, monotone convergence in
\eqref{eq:heat-exhaustion} gives
\[
 \int_0^\infty k_U(t,x,y)\dd t
  =\lim_{a\to\infty}\int_0^\infty k_U^{(-a,a)}(t,x,y)\dd t
  =\lim_{a\to\infty}G_a(x,y),
\]
with infinity allowed. This proves Lemma~\ref{lem:green}.

\section*{AI-use statement}
OpenAI's Codex assisted in developing and writing the argument, checking
cited passages, and conducting internal AI reviews of the proof and
manuscript. The mathematical claims are submitted for independent review.

\begin{thebibliography}{9}
\bibitem{FKP}
M.~Fraas, D.~Krej\v{c}i\v{r}\'ik, and Y.~Pinchover,
\emph{On some Strong Ratio Limit Theorems for Heat Kernels},
arXiv:0912.4337v3 (2010).
\url{https://arxiv.org/abs/0912.4337v3}.

\bibitem{Krejcirik}
D.~Krej\v{c}i\v{r}\'ik,
\emph{Large-time behavior of the heat equation: subcriticality versus criticality},
AIM problem contribution (2015).
\url{https://nsa.fjfi.cvut.cz/problems/04_AIM_2015/2015-18/2015-18_AIM_Krejcirik_2.pdf}.

\bibitem{AIM217}
MColbrook/AIM, \emph{217. Does subcritical heat flow always decay faster?},
problem statement at
\href{https://github.com/MColbrook/AIM/blob/a153f469ee0ea64e876b2ad97c79d252c19a3e22/problems/217-critical-subcritical-heat-ratio.md}{revision \texttt{a153f469}}.

\bibitem{Simon}
B.~Simon, \emph{A Feynman-Kac Formula for Unbounded Semigroups},
arXiv:math-ph/9907022v1 (1999).
\url{https://arxiv.org/abs/math-ph/9907022v1}.

\bibitem{Teschl}
G.~Teschl, \emph{Mathematical Methods in Quantum Mechanics: With Applications
to Schr\"odinger Operators}, Graduate Studies in Mathematics 99,
American Mathematical Society, online version of February 12,
2009. \url{https://www.mat.univie.ac.at/~gerald/ftp/book-schroe/schroe.pdf}.

\bibitem{Lalley}
S.~Lalley, \emph{Lecture 5: Brownian Motion}, online lecture notes,
Proposition 4.
\url{https://www.stat.uchicago.edu/~lalley/Courses/390/Lecture5.pdf}.
\end{thebibliography}
\end{document}
